\documentclass[11pt,a4paper]{article}

\usepackage[T1]{fontenc}
\usepackage[utf8]{inputenc}
\usepackage[brazil]{babel}
\usepackage{lmodern}
\usepackage{microtype}
\usepackage[margin=2cm]{geometry}
\usepackage{hyperref}
\usepackage{titlesec}
\usepackage{xcolor}

\hypersetup{
  colorlinks=true,
  linkcolor=black,
  urlcolor=black,
  citecolor=black,
  pdfborder={0 0 0}
}

\definecolor{sectioncolor}{HTML}{111111}
\definecolor{accentcolor}{HTML}{1A1A1A}

\pagestyle{empty}
\setlength{\parindent}{0pt}
\setlength{\parskip}{5pt}
\linespread{1.05}
\selectfont

\titleformat{\section}
  {\bfseries\large\color{sectioncolor}}
  {}{0pt}{}[\vspace{-2pt}{\color{sectioncolor}\titlerule[0.8pt]}]
\titlespacing*{\section}{0pt}{12pt}{8pt}

\newcommand{\cvrole}[4]{%
  \vspace{4pt}%
  {\bfseries\color{accentcolor}#1}\hfill{\itshape #2}\\[-2pt]%
  {\itshape #3}\hfill #4\\[6pt]%
}

\begin{document}

\begin{center}
  {\fontsize{20}{24}\selectfont\bfseries Kamille Hartmann Maccagnan}\\[8pt]
  {\large Analista de Projetos \textbar{} PMO \textbar{} Inovação e Indicadores}\\[10pt]
  Campo Comprido, Curitiba-PR \enspace\textbullet\enspace (41) 9 9928-4424 \enspace\textbullet\enspace
  \href{mailto:kamillehartmannm@gmail.com}{kamillehartmannm@gmail.com}\\[3pt]
  \href{https://linkedin.com/in/kamille-hartmann-maccagnan/}{linkedin.com/in/kamille-hartmann-maccagnan/}
\end{center}

\vspace{6pt}

\section{Resumo Profissional}

Profissional com experiência em apoio à gestão de projetos, programas e iniciativas corporativas em ambientes multinacionais e de tecnologia. Atua na organização de demandas e iniciativas, acompanhamento de cronogramas, planos de ação e indicadores, elaboração de relatórios e apresentações executivas e suporte à governança de projetos. Vivência em documentação de processos, registro de reuniões, padronização de informações e interface com stakeholders de diferentes áreas. Conhecimento em metodologias ágeis e Boas Práticas PMOBOK, com domínio de Excel, PowerPoint e ferramentas de análise de dados. Perfil organizado, analítico e comunicativo, com capacidade de transformar informações complexas em materiais claros e objetivos para apoio à tomada de decisão.

\section{Experiência Profissional}

\cvrole{Anjuss}{Analista de Suporte de TI}{07/2025 -- Atual}

Atuo na organização e priorização de demandas operacionais, com criação e atualização de documentações e POPs que asseguram padronização, rastreabilidade e governança dos processos da rede de lojas. Desenvolvo relatórios e dashboards em Power BI para acompanhamento de indicadores e desempenho, consolidando resultados e apoiando a identificação de pendências e oportunidades de melhoria. Conduzo treinamentos e ações de comunicação interna para alinhamento de equipes e divulgação de procedimentos. Participo de reuniões de acompanhamento, registrando encaminhamentos e atuando como ponte entre usuários de negócio e áreas técnicas.

\cvrole{HORSE do Brasil}{Estagiária de TI (IS/IT LATAM)}{07/2024 -- 07/2025}

Atuei no apoio a projetos e PMO em ambiente industrial multinacional, com acompanhamento de portfólio de projetos, cronogramas, prazos, entregas e responsáveis por iniciativas estratégicas. Elaborei relatórios executivos e dashboards em Power BI para consolidação de indicadores e apoio a reuniões de status com lideranças e stakeholders internacionais. Participei da organização de documentação, levantamento de riscos, análise de pendências e governança de projetos, com registro de alinhamentos e encaminhamentos em reuniões multilíngues. Utilizei o ServiceNow para gestão de demandas e atuei como ponte entre áreas de negócio, operações e fornecedores.

\cvrole{Restaurante Familiar}{Analista de Dados e Suporte Técnico}{01/2023 -- 07/2024}

Atuei em atividades analíticas e administrativas para apoio à gestão do negócio, com elaboração de relatórios, indicadores de desempenho e dashboards em Power BI para acompanhamento de resultados operacionais e financeiros. Organizei e padronizei informações e bases de controle em Excel, contribuindo para a confiabilidade dos dados e o cumprimento de rotinas e prazos. Atuei no alinhamento entre áreas operacionais e administrativas, apoiando o controle de planos de ação e a melhoria da eficiência dos processos.

\section{Competências Técnicas}

\textbf{Gestão de Projetos e PMO:} Apoio à gestão de projetos e programas, acompanhamento de cronogramas e planos de ação, controle de prazos e entregas, portfólio de projetos, governança, reuniões de acompanhamento, elaboração de pautas e registro de encaminhamentos, Boas Práticas PMOBOK, Scrum, Kanban.

\textbf{Apresentações e Comunicação:} Apresentações executivas em PowerPoint, relatórios gerenciais, transformação de informações complexas em materiais claros e objetivos, comunicação interna, interface com stakeholders, trabalho em equipe.

\textbf{Análise e Organização:} Acompanhamento de indicadores, consolidação de resultados, Excel avançado, Power BI, análise e organização de informações, documentação de processos, POPs, padronização e bases de conhecimento.

\textbf{Ferramentas:} ServiceNow, Trello, Monday, Google Sheets, Pacote Office.

\textbf{Idiomas:} Inglês Avançado \enspace\textbullet\enspace Espanhol Básico \enspace\textbullet\enspace Coreano Básico

\section{Projetos e Conquistas}

1º lugar no Transformation Day Renault 2023 com solução orientada a inovação, dados e melhoria de processos. Participação no Hackathon Bosch 2023 com foco em resolução de problemas de negócio e prototipação de soluções. Acompanhamento de portfólio de projetos e elaboração de relatórios executivos em ambiente multinacional na HORSE. Estruturação de KPIs e dashboards para monitoramento de performance. Certificação em Scrum Ágil (LinkedIn).

\section{Formação Acadêmica}

\textbf{PUCPR} \hfill Curitiba, PR\\
Graduação em Gestão da Tecnologia da Informação \hfill Previsão de conclusão: 06/2026\\[6pt]

\textbf{Escola Conquer} \hfill Curitiba, PR\\
Pós-graduação em Liderança e Gestão em Tecnologia \hfill Conclusão: 2024\\[6pt]

\textbf{Universidade Tuiuti do Paraná} \hfill Curitiba, PR\\
Graduação em Medicina Veterinária \hfill 06/2019\\[6pt]

\textbf{Qualittas} \hfill Curitiba, PR\\
Pós-graduação em Clínica Médica e Cirúrgica de Animais Exóticos e Pets Não Convencionais \hfill 12/2022

\end{document}
